% Schemi degli esercizi (usati sia nell'esercitazione sia nelle soluzioni)
\newcommand{\partitoreV}[3]{% #1 = V, #2 = R1, #3 = R2
\begin{circuitikz}[scale=0.8, transform shape, line width=0.6pt]
  \rail{(0,4.4)}{$V = #1$}
  \draw (0,4.4) -- (0,4.2) to[R, l_={$R_1 = #2$}] (0,2.2) to[short, *-] (0,2.2)
        to[R, l_={$R_2 = #3$}] (0,0.2) -- (0,0) node[sground]{};
  \draw (0,2.2) to[short, -o] (1.4,2.2) node[right]{X};
\end{circuitikz}}

\newcommand{\partitoreI}[3]{% #1 = I, #2 = R1, #3 = R2
\begin{circuitikz}[scale=0.8, transform shape, line width=0.6pt]
  \draw (0,4.4) node[ocirc]{} to[short, i>^={$I = #1$}] (0,3.2);
  \draw (-1.6,3.2) -- (1.6,3.2); \draw (0,3.2) node[circ]{};
  \draw (-1.6,3.2) to[short, i>^={$I_1$}] (-1.6,2.6) to[R, l_={$R_1 = #2$}] (-1.6,0.6) -- (-1.6,0.3) -- (1.6,0.3);
  \draw (1.6,3.2) to[short, i>_={$I_2$}] (1.6,2.6) to[R, l^={$R_2 = #3$}] (1.6,0.6) -- (1.6,0.3);
  \draw (0,0.3) node[circ]{} -- (0,0.1) node[sground]{};
\end{circuitikz}}

\newcommand{\partitorePot}{%
\begin{circuitikz}[scale=0.8, transform shape, line width=0.6pt]
  \rail{(0,4.4)}{$V = 10\,\mathrm{V}$}
  \draw (0,4.4) -- (0,4.2) to[R, l_={$R_1 = 1\,\mathrm{k\Omega}$}] (0,2.2) to[short, *-] (0,2.2)
        to[vR, l_={$R_2 = 0 \div 4{,}7\,\mathrm{k\Omega}$}] (0,0.2) -- (0,0) node[sground]{};
  \draw (0,2.2) to[short, -o] (1.4,2.2) node[right]{X};
\end{circuitikz}}

\newcommand{\resistenzaSola}[2]{% #1 = V, #2 = R
\begin{circuitikz}[scale=0.8, transform shape, line width=0.6pt]
  \rail{(0,3.4)}{$V = #1$}
  \draw (0,3.4) -- (0,3.2) to[R, l_={$R = #2$}, i>^={$I$}] (0,0.8) -- (0,0.6) node[sground]{};
\end{circuitikz}}
